% =====================================================================
%  III. METHODOLOGY
%  Source: notes/experimental-protocol.md, sections 2, 3, 6, 7 and 8.
%  Where this text goes beyond the protocol, a comment says so.
% =====================================================================
\section{Methodology}
\label{sec:methodology}

\subsection{Configuration Grid}
\label{sec:grid}

The routed pool is a grid, not a ladder. It crosses four model sizes with
two reasoning modes, giving eight cells
$\cells = \{c_1, \dots, c_8\}$ (Table~\ref{tab:grid}). All cells come from
one model family, Qwen3, which offers a size ladder in which every model can
run with or without a thinking phase~\cite{qwen3}. Using one family means
that size and mode are the only things that vary between cells: tokenizer,
training recipe and chat template stay the same. Quantisation is held fixed
at 4-bit (Q4\_K\_M) for the same reason. The cost of that choice is that
results describe quantised models, a limitation we return to in
Section~\ref{sec:discussion}.

\begin{table}[!t]
\centering
\caption{The eight-cell configuration grid. Sizes on disk are for the
Q4\_K\_M weights used in all experiments.}
\label{tab:grid}
\renewcommand{\arraystretch}{1.1}
\begin{tabular}{@{}c l l r@{}}
\toprule
Cell & Model & Mode & Weights \\
\midrule
$c_1$ & Qwen3-0.6B & no-think & 0.52\,GB \\
$c_2$ & Qwen3-0.6B & think    & 0.52\,GB \\
$c_3$ & Qwen3-1.7B & no-think & 1.36\,GB \\
$c_4$ & Qwen3-1.7B & think    & 1.36\,GB \\
$c_5$ & Qwen3-4B   & no-think & 2.62\,GB \\
$c_6$ & Qwen3-4B   & think    & 2.62\,GB \\
$c_7$ & Qwen3-8B   & no-think & 5.23\,GB \\
$c_8$ & Qwen3-8B   & think    & 5.23\,GB \\
\bottomrule
\end{tabular}
\end{table}

\subsection{Cost of a Configuration}
\label{sec:cost}

We do not assume that cells are ordered by parameter count. A thinking cell
of one size can cost more than a non-thinking cell of the next size up, and
whether it does depends on the prompt. Cost is therefore measured per
prompt $q$ and per cell $c$, in three units:
\begin{align}
  \cost_{\mathrm{tok}}(c, q) &= \alpha\, n_{\mathrm{in}}(c, q)
                              + \beta\, n_{\mathrm{out}}(c, q),
  \label{eq:tokcost}\\
  \cost_{\mathrm{lat}}(c, q) &= \text{wall-clock time of the request (s)},
  \label{eq:latcost}\\
  \cost_{\mathrm{E}}(c, q)   &= \text{GPU energy of the request (J)},
  \label{eq:energycost}
\end{align}
where $n_{\mathrm{in}}$ and $n_{\mathrm{out}}$ are input and generated
token counts, the latter including any reasoning trace, and $\alpha$ and
$\beta$ are per-token prices that make the token cost comparable with
metered API billing. Energy is measured rather than estimated
(Section~\ref{sec:energy}). All definitions below apply to any of the three
units; we write $\cost(c, q)$ when the unit does not matter.

\subsection{Minimum Sufficient Configuration}
\label{sec:cstar}

Let $s(q, c) \in \{0, 1\}$ indicate whether cell $c$ produces an acceptable
answer to $q$ under the success criteria of
Section~\ref{sec:labelling}. The \emph{minimum sufficient configuration}
is the cheapest cell that succeeds:
\begin{equation}
  \cstar(q) = \argmin_{c \in \cells \,:\, s(q, c) = 1} \cost(c, q).
  \label{eq:cstar}
\end{equation}
It is undefined when no cell succeeds. Those prompts are kept and reported
separately as the \emph{unservable fraction}. Dropping them would flatter
every policy, since no policy can serve them and all of them pay to find
that out.

\subsection{Manual Escalation and the Escalation Tax}
\label{sec:tax}

\emph{Policies.} A manual escalation policy $\pi$ consists of a starting
cell and an ordered ladder of cells to try next. The user submits $q$ to the
starting cell; when an answer is judged a failure, the \emph{whole prompt}
is resubmitted to the next rung, and so on until an answer is accepted or
the ladder runs out. For a prompt $q$, let $a_1, \dots, a_m$ be the cells
actually attempted, with $a_m$ the cell whose answer is finally accepted.
The policy's cost is the sum over every attempt:
\begin{equation}
  \mathrm{Cost}(\pi, q) = \sum_{j=1}^{m} \cost(a_j, q).
  \label{eq:policycost}
\end{equation}
Because $n_{\mathrm{in}}$ enters \eqref{eq:tokcost} at every attempt, a
long prompt makes escalation expensive even when each failed answer is
short. That resubmission cost is what users notice in practice, and it is
absent from cascade evaluations that charge only for generated tokens.

\emph{The tax.} The escalation tax is the excess over the minimum
sufficient configuration:
\begin{equation}
  \Tax(\pi, q) = \mathrm{Cost}(\pi, q) - \cost\bigl(\cstar(q), q\bigr).
  \label{eq:tax}
\end{equation}
Splitting the sum in \eqref{eq:policycost} at the final attempt gives an
exact decomposition into two terms with different causes:
\begin{align}
  \Tax(\pi, q) &= \Esc(\pi, q) + \Stk(\pi, q), \label{eq:decomp}\\
  \Esc(\pi, q) &= \sum_{j=1}^{m-1} \cost(a_j, q), \nonumber\\
  \Stk(\pi, q) &= \cost(a_m, q) - \cost\bigl(\cstar(q), q\bigr). \nonumber
\end{align}
The \emph{escalation} component is the cost of attempts that were thrown
away. The \emph{stickiness} component is what the user overpays for the
configuration that finally answered. When failures are always detected,
$a_m$ succeeded, so by \eqref{eq:cstar} its cost is at least that of
$\cstar(q)$ and $\Stk \geq 0$; both components are then non-negative. A
user who begins on the largest thinking cell and never needs to move pays
only stickiness. A user who climbs from the smallest cell and happens to
stop exactly at $\cstar(q)$ pays only escalation. Most trajectories pay
some of each, and the two respond to different remedies, which is why we
report them separately.
% NOTE for reviewers: the protocol (S3) describes stickiness as the premium
% paid when the user starts above c* and never moves. The definition above
% generalises it to the overshoot of the final cell, which also covers a
% coarse ladder that jumps past c* after escalating. The two agree for the
% sticky-high policy. Decomposition is exact either way.

\begin{table}[!t]
\centering
\caption{Manual escalation policies. Cells as in Table~\ref{tab:grid}.}
\label{tab:policies}
\renewcommand{\arraystretch}{1.1}
\small
\begin{tabularx}{\columnwidth}{@{}l >{\raggedright\arraybackslash}p{2.9cm} >{\raggedright\arraybackslash}X@{}}
\toprule
Policy & Ladder & Behaviour modelled \\
\midrule
P1 & $c_1{\to}c_3{\to}c_5{\to}c_7{\to}c_8$ & \emph{Frugal climb}: cost-conscious; larger model first, thinking last \\
P2 & $c_1{\to}c_2{\to}c_4{\to}c_6{\to}c_8$ & \emph{Think first}: turns on thinking before moving up a size \\
P3 & $c_8$ only & \emph{Sticky high}: heavyweight left selected from earlier work \\
P4 & Start drawn from an assumed distribution over cells, then the P1 ladder & \emph{Sticky mixed}: a population of users \\
\bottomrule
\end{tabularx}
\end{table}

\emph{Policy family.} We construct the comparator ourselves, so we
construct it generously. Table~\ref{tab:policies} lists four policies that
range from frugal to wasteful. P4 models a population of users who start
wherever they left off. We have no usage logs from which to estimate where
that is, so its starting distribution is an assumption: uniform over the
eight cells in the primary analysis, and, as a sensitivity analysis, a
distribution in which half of all queries start at $c_8$ and the rest are
spread evenly over the other seven cells. Both are labelled as assumed
wherever they are reported.

\emph{Failure detection.} An escalating user has to notice that an answer
is wrong. We evaluate two regimes. Under \emph{oracle detection} the user
always notices, and never discards a correct answer. This regime favours
the manual policies, so it gives a lower bound on the tax, and it is the
regime we quote in headline results. Under \emph{imperfect detection} the
user accepts a wrong answer with probability $p_{\mathrm{acc}}$ and
rejects a correct one with probability $p_{\mathrm{rej}}$. The answer from
the last rung is kept whatever its quality, because the ladder has run out.
A prompt may then end unsolved, so costs are also reported per
\emph{solved} prompt,
\begin{equation}
  \bar{\cost}_{\mathrm{solved}}(\pi) =
  \frac{\mathbb{E}_q\bigl[\mathrm{Cost}(\pi, q)\bigr]}
       {\Pr_q\bigl[\text{$q$ solved under $\pi$}\bigr]},
  \label{eq:persolved}
\end{equation}
We have no measurement of the two probabilities for users of LLM
services, so they are treated as assumptions and swept:
$p_{\mathrm{acc}} \in \{0, 0.05, 0.10, 0.15, 0.20, 0.30\}$ and
$p_{\mathrm{rej}} \in \{0, 0.05, 0.10\}$. The point $p_{\mathrm{acc}} =
0.15$, $p_{\mathrm{rej}} = 0.05$ is shown only as an illustration. Every
outcome of a policy, including its tax and its probability of solving a
prompt, is an expectation over the user's accept-or-reject decisions, so we
compute it exactly from the labelled success vectors rather than by
simulating users.

\subsection{Routers as Per-Cell Success Predictors}
\label{sec:routers}

Every router in this study predicts a vector of success probabilities, one
per cell, rather than a single class label:
\begin{equation}
  f(q) = \bigl(\hat{p}_1(q), \dots, \hat{p}_8(q)\bigr), \qquad
  \hat{p}_i(q) \approx \Pr\bigl[s(q, c_i) = 1\bigr].
  \label{eq:router}
\end{equation}
A separate rule turns the vector into a decision. Given a quality threshold
$\tau \in [0, 1]$, the router picks the cheapest cell it believes will
succeed:
\begin{equation}
  \hat{c}_\tau(q) =
  \begin{cases}
    \displaystyle\argmin_{i \,:\, \hat{p}_i(q) \geq \tau} \bar{\cost}(c_i)
      & \text{if some } \hat{p}_i(q) \geq \tau,\\[1.2ex]
    \displaystyle\argmax_{i} \hat{p}_i(q) & \text{otherwise},
  \end{cases}
  \label{eq:selection}
\end{equation}
where $\bar{\cost}(c_i)$ is the mean cost of cell $c_i$ on the training
split, because the true per-prompt cost is unknown before generation.
% ASSUMPTION beyond the protocol: the fallback when no cell clears tau
% (argmax of p-hat) and the use of training-split mean cost. Flag for review.
Separating prediction from decision has two benefits. One trained router
yields a whole cost--quality curve by sweeping $\tau$, much as Hybrid LLM
exposes a quality knob at test time~\cite{ding2024hybridllm}. And the
asymmetry between the two kinds of error can be handled by moving $\tau$
instead of retraining. Routing one cell too low produces a failure and a
retry; routing one cell too high only wastes compute. We expect these costs
to differ by a large factor, which is why accuracy against $\cstar$ is not
used as the objective. Direct eight-way classification of $\cstar(q)$ is
kept as an ablation.

\emph{What happens after a misroute.} A router that picks a failing cell
does not end the interaction; the user still needs an answer. We charge a
routed policy the same way as a manual one: the routed cell is the first
attempt, and on failure the prompt is resubmitted once to the largest
thinking cell $c_8$. As a sensitivity analysis, a failed routed attempt
instead continues up the P1 ladder from the routed cell. The escalation tax
of a router is then computed with \eqref{eq:tax} exactly as for $\pi$. For
each cost unit, recovery rule and detection regime, $\tau$ is chosen on the
validation split to minimise the net cost per solved prompt, router cost
included, and the test split is evaluated once at that value.

We compare three router families, chosen to span four orders of magnitude
in their own cost.

\emph{R1: classical classifiers.} (a) Logistic regression on TF-IDF word and
character $n$-grams. (b) Gradient-boosted trees~\cite{chen2016xgboost,ke2017lightgbm} on hand-built features: prompt length,
digit and operator density, presence of code fences, question type and
readability. (c) Logistic regression on frozen Sentence-BERT
embeddings~\cite{reimers2019sbert}. All run on a CPU.

\emph{R2: fine-tuned encoders.} DistilBERT~\cite{sanh2019distilbert} and
DeBERTa-v3-small~\cite{he2021deberta} with eight sigmoid output heads,
trained with class-weighted binary cross-entropy.

\emph{R3: a language model as router.} Qwen3-0.6B, the smallest member of
the pool, is shown the eight configurations, labelled A to H in grid
order, and asked which is the cheapest that will answer the request
correctly. It is prompted zero-shot and few-shot; the few-shot prompt adds
one training prompt for each configuration that is the $\cstar$ of some
training prompt. Rather than decode a letter, we read the model's
next-token distribution at the first answer token, keep the mass on the
eight letters and renormalise it. This constrains the answer to a valid
configuration, as constrained decoding would, and yields a distribution
over $\cstar(q)$ instead of a single guess. Treating success as monotone
in mean cost, we set $\hat{p}_i(q)$ in \eqref{eq:router} to the mass on
configurations whose mean cost does not exceed that of $c_i$. A LoRA-tuned
variant~\cite{hu2022lora} is planned
\tbd{needs the full-precision weights, not yet downloaded}. We
deliberately use the smallest pool member: a router larger than the models
it routes between defeats its own purpose.

\emph{Cost-sensitive training.} Because misrouting costs are asymmetric, a
variant of each trainable router weights errors by their realised cost
consequence, computed from the labelled per-cell costs, instead of treating
all errors equally.

\subsection{Net-of-Router Accounting}
\label{sec:net}

A router is itself a workload. Each router's latency and energy per
decision are measured with the same instruments as the pool
(Section~\ref{sec:energy}) and added to the cost of every query it routes.
For a router $r$ and a baseline policy $b$, the \emph{net saving} is
\begin{equation}
  S_{\mathrm{net}}(r; b) =
    \mathbb{E}_q\bigl[\mathrm{Cost}(b, q)\bigr]
  - \mathbb{E}_q\bigl[\mathrm{Cost}(r, q) + \cost_{\mathrm{route}}(r, q)\bigr],
  \label{eq:netsaving}
\end{equation}
where $\cost_{\mathrm{route}}(r, q)$ is the router's own cost for $q$.
Latency is measured directly, one decision at a time, as a deployed router
would make them. For token cost the term is non-zero only for R3, which
consumes input and output tokens of its own. R3's decisions run through the
labelling harness as short generations, so their tokens, latency and energy
are metered per prompt exactly as the pool's are. NVML cannot see a CPU
router, so its energy
per decision is bounded above by its decision time multiplied by the upper
end of the processor's configurable power range (54\,W), and that bound is
what we charge.
Ranking router families by gross saving and by $S_{\mathrm{net}}$ need not
give the same order, and testing whether it does is one of the aims of the
study.

\subsection{Reactive Cascade Baseline}
\label{sec:cascade}

The strongest published alternative to predicting in advance is to try the
cheap configuration and check. We implement a cascade in the style of
AutoMix~\cite{aggarwal2024automix}: generate on the cheapest cell, ask the
same model to verify its answer with a few-shot self-verification prompt,
and escalate along the P1 ladder when verification fails. The prompt's
four worked examples are generic, and none is drawn from the corpus. The
verifier's confidence is read from a single forward pass as
$P(\text{Yes}) / \bigl(P(\text{Yes}) + P(\text{No})\bigr)$ at its first
answer token, and a query escalates when that confidence falls below a
threshold chosen on the validation split, as $\tau$ is for a router. The
answer from $c_8$, the top rung, is kept without verification. AutoMix's
POMDP meta-verifier is omitted. That simplification may weaken the cascade
relative to the original method, and we keep it in view wherever the
cascade is compared. Verification calls are metered and charged like any
other generation: the check on the answer the cascade returns is part of
that attempt, and the answers and checks it passes over are escalation.
The returned answer then meets the user's check and recovery as a routed
answer does. This baseline pays for failed attempts,
as a user does, but decides with information from the answer itself.
Beating it therefore requires that predicting before generation save more
than the extra information from generating first is worth.

\subsection{Energy Measurement}
\label{sec:energy}

GPU board power is read through NVIDIA's management library, NVML~\cite{nvidia_nvml}, by a sampling thread that polls
instantaneous power every 100\,ms. Two properties of the hardware, both
established by the telemetry gate in Section~\ref{sec:gate-telemetry},
shape the method. First, the cumulative energy counter that NVML also
exposes is unusable on this GPU, so energy is always integrated from sampled
power. Second, reported power refreshes only every 0.5\,s and is smoothed.
A short no-think answer from the smallest model can finish inside a single
refresh period, so its energy cannot be read from its own time window.

Energy is therefore integrated over \emph{batch windows}. Prompts are taken
in batches of 20, and each batch is run twice back to back on the loaded
model, once per mode, so that every window holds a single cell. The mode
that runs first alternates from one batch to the next. Each window is
extended by a 6\,s tail after its last generation, because a pilot run
showed that reported power takes about 5\,s to decay once generation stops;
a shorter tail would charge one window's decay to the next. The integral is
computed with the trapezoidal rule. Readings above 200\,W, which the GPU
cannot draw under its 115\,W limit, are discarded as read errors; one such
reading appeared in the pilot run. We report two quantities: board energy,
the raw integral, and dynamic energy, which subtracts the idle power of the
GPU with the model loaded. That baseline is the mean power over 60\,s after
warm-up and 30\,s at the end of each model's session. Dynamic energy is the
figure we attribute to queries. Each query receives a share of its window's
dynamic energy in proportion to its share of the window's generation wall
time, and we check that apportioning by token share gives the same result.
A window is discarded and rerun, up to four times, if utilisation at the
end of its tail shows another workload or if any of its generations had to
reload the model. Routers that run on the GPU are measured the same way.

Two derived measures connect energy to the escalation tax. \emph{Joules per
solved query} divides the energy a policy spends, including failed attempts,
by the number of prompts it solves, so a cheap configuration that fails and
triggers a retry is charged for the retry. \emph{Carbon} is reported as
$\mathrm{gCO_2e} = (E / 3.6 \times 10^{6})\,I$ for a grid intensity $I$ in
gCO$_2$e/kWh, over a range of regional intensities rather than a single
value, following the accounting of Lacoste \emph{et
al.}~\cite{lacoste2019quantifying} and Luccioni \emph{et
al.}~\cite{luccioni2024powerhungry}. NVML reports whole-board power, so any
other GPU workload contaminates the reading; measurement windows require an
otherwise idle GPU, and utilisation is logged throughout to flag windows where
it does not match the scheduled work.
